\emph{equilibrium} kernel witnesses carry the full $1/200$
mixing at zero entropy production --- no non-equilibrium state
is needed, and whether the deployed \textsc{pure-i} wiring is
compiler-forced became the registered successor contract
(\veri{v908\_seam\_equilibrium\_wiring.py},
\texttt{SEAM.STATE.WIRING.SELECTOR.01}); the \mO\ premise stays
unmoved.  \emph{Dated addendum (2026-08-12, rounds 69--71):}
the successor contract is \emph{closed}, and the answer is a
\emph{freedom theorem}
(\veri{v911\_wiring\_freedom.py}): the exact Gr\"obner census of
the orbit/edge rules leaves exactly \emph{one} admissible
three-dimensional wiring component modulo the rule gauge (the
rotation cell); $Z$/$X$ wirings are excluded gauge-invariantly
by orientation propagation; the deployed \textsc{pure-i} point
is an \emph{interior} point of that cell, connected to the
pure-$J$ ray by an integer rule-gauge element --- \textsc{pure-i
is a deployment representative, not a compiler theorem}.  No
compiler-side demand distinguishes the $\theta_S$ sheet-swap
frame within the rule-gauge algebra (five demand classes, all
silent on the angle), the strict-collar RP Grams transform
\emph{exactly covariantly} under the rule gauge, and the same
deployed parent that fails strict collar RP at $\theta_S$
\emph{passes it exactly} in the integer-conjugated frame: the
rounds-58/59 obstruction is frame-relative (a statement at the
fixed $\theta_S$ frame, not a wiring no-go), and a canonical
wiring selection would require an additionally \emph{derived}
physical frame.  The \mO\ premise of the Wick-compiler demand
stays unmoved.  No marker moves.
\end{keybox}

\begin{keybox}[The current reduction state of P2 (2026-07-22) --- text consolidation, markers unchanged]
The axiom P2 stays \emph{declared} (\texttt{AX.P2.01}); what has changed is how much of it is now
carried by theorems. The honest ledger state: P1 ($\cthree=1/(8\pi)$) and the \emph{four marks} are
\mE-anchored (Gauss--Bonnet $n=2\chi=4$, \veri{v216\_marks\_gauss\_bonnet.py}); given the four marks,
$\gcar=5$ follows from the partition $4=2{+}1{+}1$ plus the Deligne/Birkhoff--Grothendieck/stability
typing (\veri{v491\_p2\_partition\_corollary.py}, \veri{v499\_p2\_weights\_deligne\_bg.py}; residual
R1, the module identification, \mC); and the clock's \emph{order}, the $\Z_8$ tower, the deck's
\emph{class} (a half-period translation) and the \emph{position} half of the alignment bit are
theorems (\veri{v506\_seam\_clock\_rigidity.py}, \veri{v507\_seam\_bit\_origin.py},
\veri{v510\_seam\_bit\_freedom.py}). What remains as carrier input is \emph{one discrete
symmetry-lift bit}: flag transitivity of the four marks (marks \emph{and} their two sides
indistinguishable, $V_4\to D_4$) $\iff$ $\tau=i$ (equivalently: the clock / the CM-fixed
half-period / the non-split extension class on the free seam circle) --- on the free seam circle
bare mark-transitivity is automatic for \emph{every} configuration (the pair-exchanging $V_4$) and
no local jet sees the side bit, so the flag form is the sharp discrete statement of the datum and
its negation is concretely measurable (odd-doublet split $(2/\pi)\ln\cot(\delta/2)\neq0$)
(\veri{v512\_seam\_tau\_flag.py}, \texttt{SEAM.TAU.FLAG.01}). The bit has since survived
\emph{ten} side-blind derivation attacks on that scoreboard --- free RP/$\Theta$ existence (the
eighth, \vref{v521}), the exhausted Wick-computable free-plus-twist state class (the ninth,
\vref{v525}) and the first genuinely interacting seam toy (the tenth, \vref{v529}) --- and is now
physically \emph{defined} as the \emph{twist-class choice}: facets \#14/\#15 extend the web to
$15$ exact equivalences (both directions each), the flip-axis fraction is a gauge-robust order
parameter ($O=\tfrac12$ at $m=4$, $0$ elsewhere), and its $\eta$ holonomy is in principle
interferometrically readable (\vref{v528}, \texttt{SEAM.BIT.TWISTCLASS.01}; the bit remains
formal input --- a definition replaces no derivation, no marker moves); and it now carries the
first \emph{positive} selection datum: reflection positivity, imposed on the interacting mark
family, keeps exactly \emph{one} member alive --- $\delta=\pi/2$ with positive coupling, the
carrier of that order parameter (\veri{v534\_seam\_straddle\_cone.py},
\texttt{SEAM.STRADDLE.CONE.01}; toy-level evidence for a dynamical origin, not a derivation);
plus the \mC\
identifications R1/R2
of \vref{v499}. In one sentence: \textbf{the compiler's dimensionless inputs reduce to four marks
(derived from P1-side topology), one discrete symmetry-lift bit --- the twist-class choice (flag
transitivity $\iff\tau=i$) --- and $\pi$.} This is a reduction \emph{state}, not an abolition of
the axiom: \texttt{AX.P2.01} remains the declared interface input, and its dynamical realisation
(\texttt{QGEO.REALIZE.01}) stays \mC/\mO.
\emph{[2026-08-07 addendum:} the finite normal form of the interface is now assembled as
\emph{one} object --- $C_{\rm fin}=(V,\hbar,q^{*},\sigma,\iota)$, the keybox above
(\veri{v845\_cfin\_normal\_form.py}) --- with the code-to-matter kill built in as a checked
negative; its honesty section types exactly what the assembly does \emph{not} remove: P2's
remaining content is the \texttt{v799} residual R1$'$ (the three seam axioms, the OPE-linearity
scope, the \texttt{v776} uniqueness selection).  The assembly narrows; it does not eliminate.
Markers unchanged.]
\end{keybox}

\begin{keybox}[Seed normal form, the factorial spine, and the degree-$2$ pattern \mE/\mC\ (\veri{v106\_review\_validation.py})]
Three review-validated sharpenings of the anchor picture. \textbf{(1) Seed normal form \mE.} The seed
formula is exactly
\[
  \boxed{\ \phiz=\frac{|\mu_4|}{\Nfam}\,\cthree\;+\;\Oadm\,\cthree^{\,|\mu_4|}\ }
  \qquad\Bigl(\tfrac43=\tfrac{|\mu_4|}{\Nfam},\ \ 48=\Oadm=\Nfam\dim S^+\Bigr),
\]
a \emph{linear seam term} plus a \emph{topological fourth-order correction} over all admissible fermion
states --- a structured anchor readout, not a fitted expression. \textbf{(2) Factorial spine \mE.} With
$X=6Y$ on one $16$-state generation, $\operatorname{Tr}X=0$ and $\operatorname{Tr}X^2=120=5!$ (computed
from the explicit charge table); the \emph{same} factorial moment is $|R^+(E_8)|=\sum(\text{$E_8$
exponents})=120$, doubles to $|R(E_8)|=240$, and scales to $|W(D_5)|=2^4\cdot5!=1920$ (the Hawking-power
fingerprint) --- one moment, four projections (hypercharge, roots, Weyl group, horizon). \textbf{(3) The
degree-$2$ inventory.} The compiler repeatedly selects \emph{degree two}: the Pascal closure
$2^{g-1}=\sum_{k\le2}\binom gk$ truncates at $K{=}2$ with the \emph{unique} solution $g=5$ (scan
$1..40$); the glue norms add to the root norm, $\tfrac54+\tfrac34=2$; $\cthree=\tfrac12\times\tfrac1{4\pi}$
(variational half $\times$ two-dimensional-slice Gauss--Bonnet); the pair sector $\binom52=10=\AL$; the
hypercharge moment is a \emph{second} moment; the seam is codimension \emph{two}. \emph{Named hypothesis
\mC\ (recorded, not claimed):} \textbf{Quadratic Boundary Locality} --- the seam supports only bilinear
data, so the truncation $K=2$ is \emph{forced} rather than chosen; if proven, the Pascal \emph{selection}
(red-team Target~B residual, typed \mO/\mC) upgrades to \mE\ and the chain locality $\Rightarrow K{=}2
\Rightarrow g{=}5\Rightarrow16\Rightarrow D_5\Rightarrow E_8$ closes the carrier choice. \emph{Audit:}
the even-code transition reading $240=16\times15$ is exact arithmetic but non-unique ($240=16\times5\times3$
is the established glue factorisation) and stays recorded, not promoted.
\smallskip\\
\noindent\textbf{The Sheet-Pairing Lemma (\veri{v109\_sheet\_pairing.py}) --- the first exact anchor of
the hypothesis.} Character-exact weight-multiset combinatorics of the $D_5$ spinor weights
$(\pm\tfrac12)^5$ gives: \emph{(i)} \textbf{no scalar within a sheet} --- the zero-weight multiplicity of
$S^+\!\otimes S^+$ is $0$, because the odd slot count $g=5$ flips the sign parity of $-w$ (a parity
theorem of the five-slot code); \emph{(ii)} \textbf{the scalar lives across the sheets}:
\[
  \boxed{\ S^+\!\otimes S^-=\Lambda^0\oplus\Lambda^2\oplus\Lambda^4\ }\quad(256=1+45+210),
\]
exact as weight multisets, with zero-mode grading $(1,5,10)$ $=$ the carrier code; \emph{(iii)}
\textbf{sheet-diagonal $=$ odd forms}: $(S^+\!\otimes S^+)\cup(S^-\!\otimes S^-)=2\Lambda^1+2\Lambda^3+\Lambda^5$
exactly --- within-sheet pairings are odd-form-valued, never scalar; \emph{(iv)} the scalar-bearing tower
\textbf{tops at $\Lambda^{2K}$} with $K=2$ --- bilinear seam data reach exactly the Pascal-ladder degree
(the source-layer transport theorem; document~3). \emph{Consequence \mC:} the seam's single scalar two-point kernel is necessarily a
\textbf{sheet pairing} (chirality-off-diagonal --- the same $\Z_2$ as the glue chiralities and the
branch-kernel selection); the remaining QBL step is to show the Calder\'on kernel supplies exactly this
pairing and nothing beyond.
\smallskip\\
\noindent\textbf{The Calder\'on-sheet selection theorem (\veri{v110\_calderon\_sheet.py}) --- which
involutions certify a scalar kernel.} On $H=S^+\!\oplus S^-$ the certified bilinear data of a Calder\'on
involution $\varepsilon$ ($\varepsilon^2=1$) are the matrix elements $\langle f,\varepsilon g\rangle$.
The lemma's singlet counts per sheet block, $(0,0,1,1)$, give the exact selection statement \mE:
\[
  \boxed{\ \text{a scalar two-point datum exists}\iff\varepsilon\ \text{is sheet-odd}\ }
\]
--- a sheet-odd involution (swapping $S^+\!\leftrightarrow S^-$) certifies exactly $1+1=2=|\Z_2|$
invariant scalar kernels (one per orientation), a sheet-even involution certifies \emph{none}. The
kernel count $2=|\Z_2|$ matches exactly the \emph{two Lagrangian glues} (swapped by single spinor
conjugation): the Calder\'on scalar channel carries precisely the glue ambiguity, resolved by the same
sheet choice. \emph{Anti-overclaim control \mE:} the parity theorem is generic over the odd ladder
($g=3,5,7$: within-sheet count $0$, cross count $2^{g-1}$) --- sheet-oddness forces the
\emph{half-spinor relation} $K=(g-1)/2$, \textbf{not} the value $g=5$; the $g$-selection stays with the
independent rank-$8$/integrality closure. \emph{Consequence:} the QBL residue splits into two named
parts \mC\ (recorded, not claimed): \textbf{(a) interface} --- ``the seam Calder\'on involution is
sheet-odd'' (structurally natural: the one-sided collar/double-cover deck is the \emph{same} $|\Z_2|$
that halves $\cthree=\tfrac12(4\pi)^{-1}$, and exchanging cover sides is sheet-odd by construction);
\textbf{(b) analytic core} --- ``the two-point kernel certifies only slot-degrees $\le2$''. Proving
both closes QBL${}\Rightarrow K{=}2\Leftrightarrow g{=}5$.
\smallskip\\
\noindent\textbf{The Quadratic-Transport Theorem (\veri{v111\_quadratic\_transport.py}) --- the
transport half of the analytic core \emph{is} a theorem.} Grade seam transport by Clifford degree in
the ten field operators on the five-slot Fock space ($S=\Lambda^\bullet(\mathbb C^5)$, $S^+$ the
$16$-state code; integer Jordan--Wigner model, spanning claims by full-rank certificates mod $p$,
valid over $\mathbb Q$). \emph{Parity} \mE: all linear words are sheet-odd, all $45$ quadratic words
sheet-even --- code-preserving transport has \emph{even} Clifford degree (the same $\Z_2$ again).
\emph{Minimality} \mE: sheet-even words of degree $\le1$ are the scalars alone. \emph{Generation and
completeness} \mE: the $45$ quadratics ($=\mathfrak{so}(10)$, exactly the $\Lambda^2$ term of the
certified tower $1+45+210$) generate the whole code from the vacuum \emph{and} span the full operator
algebra $\operatorname{End}(S^+)$ ($\dim=256$) at word length two --- every code operation whatsoever
is a word in pair transport, and all higher even-degree elements (e.g.\ the $210$ quartics) are
redundant. Hence
\[
  \boxed{\ \text{the transport degree is exactly }2\ }
\]
--- degree $\le1$ generates nothing, degree $2$ generates everything: a theorem, not a hypothesis.
\emph{Anti-overclaim control \mE:} the $g=3$ world behaves identically ($15$ quadratics generate
$\operatorname{End}(S^+)=16$) --- the theorem selects the \emph{degree}, not the rank; the
$g$-selection stays with Pascal closure $+$ rank-$8$. \emph{Residue after \vref{v109}--\vref{v111}} \mC: QBL's
remaining content is two \emph{interface} statements --- (a) the seam Calder\'on involution is
sheet-odd; (b$'$) the certified \emph{state} inventory is the $\le2$-slot tower. Nothing about
transport remains hypothetical.
\smallskip\\
\noindent\textbf{The Self-Counting Channel (\veri{v112\_selfcounting\_channel.py}) --- the Pascal
closure is an \emph{identity} of the certified channel.} For odd $g$, negation $w\mapsto-w$ maps the
$S^+$ weights bijectively into the opposite sheet, so the Cartan-neutral kernels of the certified
channel are exactly the pairs $(w,-w)$ --- \emph{one per code state}:
$\#\{\text{neutral kernels}\}=\dim S^+=2^{g-1}$ exactly \mE. Grading the \emph{same} neutral set by
pair-degree $m$ gives sizes $\binom gm$, $m=0..K$ --- for $g=5$ the triple $(1,5,10)$ \mE. Two
countings of one set:
\[
  \boxed{\ 2^{g-1}=\sum_{m\le K}\binom gm\ \text{is the channel counting itself}\ }
\]
(symbolically verified for all odd $g=3..13$) --- the Pascal \emph{closure} (\texttt{v2}/\texttt{v108}) is not an
extra requirement but an identity of the certified channel. \emph{Hodge fold} \mE: the code's
minus-sign grading $\binom g{2j}$ equals the folded binomial $\binom g{\min(2j,\,g-2j)}$ with
$\min(2j,g-2j)\le K$ --- every code sector appears at pair-degree $\le K$. \emph{Anti-overclaim:}
the identity holds for every odd $g$ and carries no $g$-selection; the selection stays with the
rank-$8$ closure $g+\Nfam=8$ and ladder integrality. \emph{Residue re-type \mC\ (recorded, not
claimed):} clause (b$'$) loses its closure burden --- given the sheet-odd kernel, the channel
\emph{automatically} contains one neutral kernel per code state at pair-degree $\le K$; what remains
physical is the single input ``the seam certifies through one scalar two-point kernel'' (for the
established free $c=8$ seam net, Wick extends the one kernel to the entire even tower by
determinants). QBL is thereby the structural consistency frame of the carrier choice; the
\emph{selection} is carried by rank-$8$ $+$ integrality.
\smallskip\\
\noindent\textbf{One kernel is the whole net (\veri{v113\_quasifree\_kernel.py}) --- the QBL input
merges with the holomorphy premise.} The remaining input is now anchored in the established Gate-A
objects. \emph{Bookkeeping} \mE: $c(D_5)_1=\tfrac{45}9=5=g_{\mathrm{car}}$, $c(A_3)_1=\tfrac{15}5=3=\Nfam$,
and the carrier net $(D_5)_1\times(A_3)_1=SO(10)_1\times SO(6)_1$ is $10+6=16$ \emph{free Majorana
fermions} ($c=\tfrac{16}2=8$); the whole extension tower carrier$\,(\mu{=}16)\to SO(16)_1(\mu{=}4)\to
(E_8)_1(\mu{=}1)$ carries the \emph{same} sixteen fermions --- only the certified glue grows (index
$2\times2=4=|\mu_4|$). \emph{Wick/Pfaffian} \mE: in the exact ten-Majorana model \emph{all} $210$
vacuum four-point and all $210$ six-point functions equal the Pfaffian of the single two-point kernel
--- one kernel determines all correlations. \emph{The kernel is a Calder\'on involution of rank $g$}
\mE: $M=\mathbf1+iA$ with $A^2=-\mathbf1$, so $P=M/2$ is a projection with $\operatorname{rank}P=5=
g_{\mathrm{car}}$ and $\varepsilon=iA$ an involution; at seam level (sixteen Majoranas) the rank is
$8=\operatorname{rank}E_8$:
\[
  \boxed{\ \text{the central charge is the rank of the one kernel}\ }
\]
($c=5$ carrier block, $c=8$ seam hull); and the polarization fixes the vacuum \emph{uniquely} (joint
annihilator kernel one-dimensional). \emph{Premise merge \mC\ (recorded, not claimed):} ``single
scalar two-point kernel'' $=$ ``the seam state is quasi-free'' $=$ the defining property of the
free-fermion net the holomorphy gate already posits --- QBL adds \emph{no} assumption beyond Gate~A;
when the gate closes, the carrier choice closes \mE\ along free net $\Rightarrow$ one kernel
$\Rightarrow$ sheet-odd $\Rightarrow$ self-counting tower $\Rightarrow$ complete transport
$\Rightarrow$ carrier. Gate typing \mC/\mO\ unchanged.
\end{keybox}

\begin{keybox}[Hardening (P1): the Gauss--Bonnet origin of $c_3=\tfrac1{8\pi}$]
$\cthree=1/(8\pi)$ is not an arbitrary number: it is the \emph{one-sided Gauss--Bonnet normalisation of
the seam}. The seam construction compactifies the oriented two-dimensional normal slice $N_\Sigma$ to a sphere $S^2$,
so Gauss--Bonnet gives the seam curvature integral
\[
  \oint_{S^2}K\,dA=2\pi\,\chi(S^2)=4\pi\qquad(\chi(S^2)=2),
\]
and the boundary datum is \emph{one-sided} (the Calder\'on/double-cover construction uses a single
collar), contributing a factor $1/|\Z_2|=\tfrac12$. Hence
\[
  \boxed{\ \cthree=\frac{1}{|\Z_2|\,\oint_{S^2}K\,dA}=\frac{1}{2\cdot 4\pi}=\frac{1}{8\pi}\ }
\]
\noindent equivalently $\cthree=\tfrac12(4\pi)^{-1}$, the one-sided $2$D heat-kernel coefficient.
This locates both factors precisely:
\begin{itemize}\setlength\itemsep{1pt}
\item the \emph{discrete} $8=|\Z_2|\cdot2\cdot\chi(S^2)=2\cdot2\cdot2$ \emph{coincides with} $h(D_5)=\operatorname{rank}E_8=\varphi(30)$
--- the geometric and the bootstrap readings of the ``$8$'' agree;
\item the \emph{continuous} $\pi$ is the irreducible Gauss--Bonnet / $2$D-Gaussian primitive ($\int e^{-x^2}=\sqrt\pi$), the single primitive the self-consistency loop also isolates.
\end{itemize}
\emph{Status:} this \emph{hardens} (P1) --- it gives $\cthree$ a named geometric origin and reconciles the
``$8$'' with $E_8$ --- but does \emph{not} eliminate $\pi$; $\pi$ remains the one continuous primitive. \mC
\end{keybox}

\par\medskip
\noindent\textbf{The one-paragraph version.}
There is a boundary (the seam). It yields $\cthree=1/(8\pi)$ and a five-slot carrier $\gcar=5$,
split $3$ colour $+\,2$ weak. Two base charges $y_-=-1/3$ and $y_+=1/2$ generate the whole
fermion family; the $16$ even-occupation states are one generation, three families give
$3\times16=48$. One seed $\useed=\varphi^{\mathrm{ret}}_0$ drives four low-energy observables. Finally
$\ainv\approx137$ is the exponential engine: the electroweak scale takes one fifth of it
($e^{-\ainv/5}$), the cosmological constant takes it twice ($e^{-2\ainv}$), and the Hubble scale is
the square root of $\Lambda$. A small grammar generates many sectors.
\par\medskip

% =====================================================================
\section{Headline: the Pascal compiler on five carrier slots}
% =====================================================================

The strongest new statement is that three of the theory's separate structures are the \emph{same}
truncated row of Pascal's triangle on the five-slot carrier.

\begin{keybox}[Family $=$ action ladder $=$ transition count (one Pascal row)]
The even exterior algebra of $E$ ($\dim E=\gcar=5$), the action ladder, and the $E_8$ root count are
one Pascal row:
\[
  S^+=\Lambda^0E\oplus\Lambda^2E\oplus\Lambda^4E,\qquad
  (\dim\Lambda^0,\dim\Lambda^2,\dim\Lambda^4)=(1,10,5),\qquad
  \dim S^+=1+10+5=16,
\]
\[
  \AEW:\AH:\AL=\binom50:\binom51:\binom52=1:5:10
  =(\dim\Lambda^0E,\ \dim\Lambda^4E,\ \dim\Lambda^2E),
\]
\[
  |R(E_8)|=\dim S^+(\dim S^+-1)=16\cdot15=240 .
\]
Family (exterior degree), cosmology (action charge) and $E_8$ (root transitions) read the same
$\{1,5,10\}$. \mE
\end{keybox}

\begin{figure}[ht]
\centering
\begin{tikzpicture}[>=Stealth]
  \node[font=\bfseries\small] at (0,1.25) {family (valencies)};
  \node[font=\bfseries\small] at (5.0,1.25) {action (multiplicities)};
  \node[dbox] (v0) at (0,0.45)  {$\dim\Lambda^0E=1$};
  \node[dbox] (v2) at (0,-0.45) {$\dim\Lambda^2E=10$};
  \node[dbox] (v4) at (0,-1.35) {$\dim\Lambda^4E=5$};
  \node[abox] (m0) at (5.0,0.45)  {$A_{\mathrm{EW}}=1$};
  \node[abox] (m1) at (5.0,-0.45) {$A_{H}=5$};
  \node[abox] (m2) at (5.0,-1.35) {$A_{\Lambda}=10$};
  \draw[flow] (v0)--(m0);
  \draw[flow] (v4)--(m1);
  \draw[flow] (v2)--(m2);
  \node[font=\scriptsize] at (2.5,-2.15)
    {even Hamming scheme on $C^+$: valencies $(1,10,5)\ \leftrightarrow\ $ spectral multiplicities $(1,5,10)$};
\end{tikzpicture}
\caption{Exterior--Action duality. The Standard-Model family is the valency side of the even Hamming
scheme on the five-slot code; the cosmological action ladder is its spectral (multiplicity) side. The
$\Lambda$ (pair) sector $\Lambda^2E$ has dimension $\binom52=10$, the origin of the decuple exponent.}
\label{fig:duality}
\end{figure}

\noindent Natural sector assignment:
\[
  \boxed{\ EW\leftrightarrow\Lambda^0E\ (\text{singlet}),\quad
  H\leftrightarrow\Lambda^4E\cong E^\vee\ (\dim5),\quad
  \Lambda\leftrightarrow\Lambda^2E\ (\dim10).\ }
\]
The cosmological constant lives on the \emph{pair} sector $\Lambda^2E$, $\dim=\binom52=10$; this is
why the decuple bridge has exponent $10$ --- $\Lambda$ is a product over all $\binom52$ carrier pairs
($K_5$ edges), not over the full $16$-state family.

\subsection{$\gcar=5$ as the Pascal closure}
A family is exactly the complete Pascal truncation to degree two:
\begin{equation}
  \boxed{\ \dim S^+=2^{\gcar-1}=\binom{\gcar}{0}+\binom{\gcar}{1}+\binom{\gcar}{2}\ \Longleftrightarrow\ \gcar=5\ }
  \qquad(16=1+5+10),\quad\mE,
\end{equation}
the unique positive-integer solution of $f(g):=2^{g-1}-1-\tfrac{g(g+1)}{2}=0$. \emph{Proof of
uniqueness.} Direct evaluation gives $f(1){=}{-1}$, $f(2){=}{-2}$, $f(3){=}{-3}$, $f(4){=}{-3}$,
$f(5){=}0$, $f(6){=}{+10}$; so $g{=}5$ is the only root in $\{1,\dots,6\}$. For $g\ge6$ the forward
difference is $f(g{+}1)-f(g)=2^{g-1}-(g{+}1)\ge 2^{5}-7=25>0$ and increasing, so $f$ is strictly
increasing on $[6,\infty)$ from $f(6){=}10>0$; hence $f(g)>0$ for all $g\ge6$ and there is \emph{no}
further root. Thus $\gcar=5$ is the unique solution (no finite-range caveat). \veri{v2\_carrier\_pascal.py}
This is sharper than $2\gcar=\binom{\gcar}{2}$: the family code (even states, $2^{g-1}$) equals the action
code (sum of the first three Pascal entries) only at $\gcar=5$.

\begin{keybox}[Carrier uniqueness theorem: $\gcar=5$ and the $3{+}2$ split are forced \mE]
The carrier rank and its colour/weak split are not chosen --- they are the unique solutions of two
integer conditions tied to the $E_8$ closure and the SM gauge algebra:
\begin{align}
  &\boxed{\ \gcar=5\ \text{is the unique }g\text{ with } \Nfam=\tfrac{2^{g-1}-1}{g}\in\Z^+\ \text{and}\ \operatorname{rank}E_8=g+\Nfam=8\ }\\[2pt]
  &\boxed{\ (b,s)=(3,2)\ \text{is the unique split of }\ b+s=\gcar=5,\ \ b^2+s^2=13=|R(A_3)|{+}1\ }
\end{align}
The integer-family carriers are $g\in\{3,5,7,\dots\}$ (rank $4,8,16$), so $\operatorname{rank}E_8=8$ picks
$\gcar=5$ uniquely; and $b^2+s^2=13$ with $b{+}s=5$ forces $(3,2)$, reproducing
$\dim\mathfrak g_{\mathrm{SM}}=(b^2{-}1)+(s^2{-}1)+1=12=|R(A_3)|$. The induced hypercharge is anomaly-free
($\operatorname{Tr}Y=\operatorname{Tr}Y^3=0$ over the $16$ states). This hardens P2: the \emph{rank} and
the \emph{split} are forced; only the seam$\to$carrier interface itself remains a declared input.
\veri{v14\_carrier\_uniqueness.py}
\end{keybox}

\par\smallskip\noindent\textbf{The Pascal grammar: one $\Lambda^\bullet$ mechanism in four places (anti-numerology) \mE.}
The recurring small integers are \emph{not} coincidences --- they are the \emph{same} exterior-algebra
(Pascal) mechanism read in four spaces, which is why the binomials of rows $4$ and $5$ keep reappearing:
\[
  \renewcommand{\arraystretch}{1.25}
  \begin{array}{l|l}
  \textbf{space} & \textbf{Pascal reading}\\\hline
  \text{carrier (one generation)} & \dim S^+=\Lambda^{\mathrm{even}}(5)=\tbinom50{+}\tbinom52{+}\tbinom54=1{+}10{+}5=16\\
  \text{scale ladder (cosmology)} & \AEW:\AH:\AL=\tbinom50:\tbinom51:\tbinom52=1:5:10\\
  \text{flavor } u/d\text{ readout} & \|\operatorname{Pl}(K)\|_1=\tbinom40{+}\tbinom41{+}\tbinom42=1{+}4{+}6=11=16{-}\gcar\\
  \text{mass volume} & \det M_{\mathrm{SM}}\sim(\phiz)^{6+9+10}=(\phiz)^{25}=(\phiz)^{\gcar^2}\\
  \end{array}
\]
So the carrier is the even row $5$, the scale ladder is the lower row $5$, the $u/d$ leg is the cumulative
row $4$ up to $\deg u^c=2$, and the mass volume closes on $\gcar^2$. \emph{One} grammar
($\Lambda^k$ of the $5$-slot carrier and the $4=|\mu_4|$ family fundamental), not four lucky hits.
\veri{v44\_carrier\_exterior.py}\,\allowbreak\veri{v45\_family\_exterior.py}\,\allowbreak\veri{v46\_grand\_mass\_volume.py}
\par\smallskip

\subsection{One master moment generates $40$, $41$, $240$, $\dtwo$}
With $X=6Y$, the quadratic hypercharge trace on one family is $\operatorname{Tr}_{S^+}X^2=120=5!$,
and it generates the whole integer chain:
\begin{equation}
\begin{gathered}
  \sum_{f,j}L_{f,j}=\frac{\operatorname{Tr}_{S^+}X^2}{3}=40,\qquad
  10b_1=\frac{\operatorname{Tr}_{S^+}X^2}{3}+1=41,\\[2pt]
  |R(E_8)|=2\operatorname{Tr}_{S^+}X^2=240,\qquad
  \dtwo=4!\,\operatorname{Tr}_{S^+}X^2\,\cthree^8 .
\end{gathered}
\end{equation}
The flavor budget, the $U(1)$ budget, the $E_8$ count and the second seam defect all hang on the
second hypercharge moment $5!$. \mE\ \emph{Round-22 sharpening:} the moment itself is one symbolic
identity, $\operatorname{Tr}_{S^+}X^2=4|X|^2=4h(E_8)=120$ --- the Coxeter number carries the whole
chain (\veri{v810\_cutcode\_moment.py}) \mE.

